\section{Background and related work}

\paragraph{Event-sequence and tabular transformers.} Encoder-only, masked-pretrained
transformers~\cite{vaswani2017attention,devlin2019bert} over tokenised tabular/event sequences
are established: BEHRT for medical records~\cite{li2020behrt} and
TabBERT/TabFormer for transactions~\cite{padhi2021tabular} introduced the hierarchical
field-then-sequence design our FFM follows, related tabular transformers add contextual and
row-attention encodings~\cite{huang2020tabtransformer,somepalli2021saint}, and later work
studied numeric encodings (piecewise-linear, periodic)~\cite{gorishniy2022embeddings}.
Industrial financial foundation models scale this recipe to billions of banking
events~\cite{braithwaite2025spending}; we discuss the two closest to our setting below. The neural
\emph{architecture} we build on is therefore not itself novel; our object of study is the
\emph{transfer} recipe (pretrain once, adapt with a light head) and, specifically, what it misses
relationally.

\paragraph{Industrial FFMs and the relational gap.} Two recent industrial systems make the FFM
recipe concrete at scale, and are the closest prior work to ours. PRAGMA~\cite{ostroukhov2026pragma}
pre-trains a masked transformer on heterogeneous, multi-source banking event sequences and serves
credit scoring, fraud, and lifetime-value prediction from one representation read by a linear probe
or a lightweight fine-tune --- the exact recipe we reproduce at laptop scale. TREASURE~\cite{yeh2025treasure}
encodes payment-network transactions with dedicated static- and dynamic-attribute modules and,
notably, ingests \emph{payment-network signals} (response codes, system flags) alongside consumer
behaviour, reporting large gains on abnormal-behaviour detection. Both establish the value of the
pretrain-once, adapt-lightly recipe on real financial data at industrial scale. Crucially, however,
each encodes a \emph{single entity's own event stream}: PRAGMA a user's banking history, TREASURE a
consumer's transactions --- and the network signals TREASURE reads are attributes \emph{recorded on
that consumer's own events}, not the concurrent activity of a shared merchant or IP across
\emph{other} consumers. This is precisely the structural blindness we target: the signal in a
compromised merchant or a fraud-farm IP lives in \emph{concurrent, cross-account} activity that no
per-entity encoder observes at inference. Our cross-sequence attention is therefore
\emph{orthogonal and composable} --- it augments rather than replaces a PRAGMA- or TREASURE-style
backbone, adding a third attention over a shared entity's events from other accounts while leaving
the per-entity encoders and the pretraining recipe intact.

\paragraph{Generative recommender foundation models.} Per-user sequence transformers underpin
sequential recommendation~\cite{kang2018sasrec,sun2019bert4rec}, and HSTU, OneRec, and
TIGER~\cite{zhai2024hstu,deng2024onerec,rajput2023tiger} reformulate
recommendation as sequential transduction over a user's interactions and scale to enormous sizes.
These are architecturally the same per-user sequence transformers, yet highly effective, for two
reasons we make precise: their pretraining objective \emph{is} the downstream task (next item) and
they train end-to-end (no proxy, no frozen bottleneck), and the relational signal they rely on
(collaborative affinities) is largely \emph{stationary} and absorbs into shared item embeddings ---
unlike the \emph{transient} cross-entity state (``compromised right now'') that characterises much
fraud, which must be read from concurrent cross-sequence activity.

\paragraph{GBDTs and graph methods for fraud.} Gradient-boosted trees~\cite{ke2017lightgbm,chen2016xgboost}
with client- and entity-level aggregates remain the production and competition standard. On graph
fraud, methods such as CARE-GNN and PC-GNN~\cite{dou2020caregnn,liu2021pcgnn} beat plain GNNs by
\emph{similarity-based neighbour selection} to
resist \emph{camouflage} (fraud nodes wiring mostly to legitimate ones); we return to this when we
find plain neighbour attention underperforms on camouflaged graphs (\S\ref{sec:boundary}). Our
contribution is orthogonal: a sequence-native cross-sequence attention for FFMs, and a map of when
it beats a scalar aggregate.
